\scanpage{71}

$n=\operatorname{ppcm}(d,d')$ (pour exprimer que $u$ est d'ordre $n$ exactement) -- $S$ se définissant en prenant l'orbite d'un $s_0\in E\setminus(D_\zeta\cup D_{\zeta'})$, \deleted{et $A$} i.e. l'ens des $s_i=u^i s_0$, \deleted{et} $(i\in\mathbb Z/n\mathbb Z)$, et $A$ est formé des droites $a_i=\operatorname{dr}\{s_i,s_{i+1}\}$ (il faut encore vérifier, ou exiger, que $a_i$ n'est incident qu'à $s_0$ et $s_1$, cette condition est laissée en exercice...

\rule{1.7cm}{.35pt}

Par contre, si on exige \inserted{de plus} qu'il existe une \deleted{\unsure{s|$\delta$}}métrique non dégénérée invariante par les $\unsure{s|u|\delta}\in G^+$, et en insistant pour que les automorphismes soient des \uline{rotations} (comme l'avait suggéré \unsure{Olivier Laborde|Olivier Leborgne}), on sauve la mise. En effet, les rotations sont unimodulaires, donc on trouve que les valeurs propres de $u$ sont $\{\zeta,\zeta^{-1}\}$, et on retrouve les cas précédents.... Je n'ai pas vérifié si ça marche encore en exigeant simplement que les $\in G^+$ respectent la métrique, sans plus.....

\rule{2.2cm}{.35pt}

% archive sheet number: 70
